% ============================================================
%  APPENDIX TO CHAPTER 2 -- EXERCISES S2.1-S2.21
%  Imported from Tao I/II, ECNU, USTC, Yu Pin.  Not Rudin's.
% ============================================================
\ssection{Appendix exercises}

\begin{roadmap}[How these differ from Rudin's thirty]
Rudin's own Exercises 2.1--2.30 end the chapter proper. They drill his
vocabulary: closures, condensation points, separability, the distance
function. These twenty-one come from the four comparison texts and do
something his cannot --- they make the machinery \emph{pay}.

\medskip
\raggedright
\textbf{Group I (1--10): the completeness theorems, proved with each other.}
Rudin derives everything from the least-upper-bound axiom in one direction.
These run the implications the other way and, in exercise 1, show what fails
over $\Q$ --- the exercise that turns the equivalence list of Chapter 1 into
a statement about $\R$ specifically.

\medskip
\textbf{Group II (11--16): metric spaces past Rudin's stopping point.}
Completion, total boundedness, the contraction mapping principle, and the
closed bounded complete set that is still not compact.

\medskip
\textbf{Group III (17--20): size.} Uncountability without a diagonal
argument, Cantor's theorem, Schr\"oder--Bernstein, and a compact space
containing a sequence with no convergent subsequence.

\medskip
\textbf{Group IV (21): topology as a tool.} Furstenberg's proof that there
are infinitely many primes, using nothing but Rudin's 2.23.

\medskip
\textbf{Marked $\star$:} harder, or requiring a result from a later chapter.
\end{roadmap}

\begin{enumerate}[label=\textbf{S2.\arabic*.}, leftmargin=3.4em, itemsep=9pt,
                  topsep=8pt, labelsep=0.5em]

% ---------------- Group I ----------------
\item Give examples showing that over the rationals, each of the following
fails: the least-upper-bound principle; the monotone convergence theorem;
the cluster-point (Bolzano--Weierstrass) theorem; and the Cauchy
criterion.\kn{The exercise that gives the equivalence list of Chapter 1 its
point. One example does all four: take rationals increasing to $\sqrt2$ ---
say the decimal truncations. The set is bounded above with no rational
supremum, the sequence is monotone bounded and does not converge in $\Q$, it
is Cauchy without a limit, and its only candidate cluster point is
irrational. That a single witness kills all four is the equivalence, seen
from below.}
\provenance{ECNU \cjk{习题}{Ex} 7.1 \#4 (bk.\ p.\,155).}

\item Let $H = \bigl\{\bigl(\tfrac{1}{n+2}, \tfrac1n\bigr) :
n = 1,2,3,\dots\bigr\}$. (a) Does $H$ cover $(0,1)$? (b) Can finitely many
members of $H$ cover $\bigl(0,\tfrac12\bigr)$? (c) Can finitely many cover
$\bigl(\tfrac{1}{100},1\bigr)$?\kn{Three questions probing exactly where
finiteness breaks. Note that (c) succeeds while (b) fails, which isolates
the mechanism: the trouble is never ``the set is infinite'' but ``the cover
degenerates at a point the set does not contain.'' Rudin's cover examples at
2.32 are one-liners with no such sub-questions.}
\provenance{ECNU \cjk{习题}{Ex} 7.1 \#5 (bk.\ p.\,155).}

\item Prove the cluster-point theorem (S2.5) using the finite covering
theorem (S2.4).\kn{The contrapositive direction Rudin never runs. If bounded
infinite $E$ had no cluster point, every $x \in [a,b]$ would have a
neighbourhood meeting $E$ in finitely many points; these cover $[a,b]$, so
finitely many do, making $E$ finite. Compactness converting ``locally
finite'' into ``globally finite'' is the whole idea, and it recurs at 4.15.}
\provenance{ECNU \cjk{习题}{Ex} 7.1 \#8 (bk.\ p.\,155).}

\item Prove the Cauchy criterion using the cluster-point theorem.\kn{A Cauchy
sequence is bounded, so it has a cluster point; then show the whole sequence
converges to it, using the Cauchy condition to transfer nearness from the
subsequence to the tail. Rudin does exactly this at 3.11, without naming the
route. Together with exercise 3 this closes the loop of implications.}
\provenance{ECNU \cjk{习题}{Ex} 7.1 \#9 (bk.\ p.\,155).}

\item $\star$~Prove the finite covering theorem using the least-upper-bound
principle, without bisection.\kn{The ``creeping along'' proof, and a genuinely
different mechanism from S2.4. Let $E$ be the set of $x \in [a,b]$ such that
$[a,x]$ admits a finite subcover, and put $c = \sup E$. Show $c \in E$ ---
some cover element contains $c$, and it also contains points slightly left of
$c$ --- and then that $c = b$, since otherwise the same element lets you
creep past $c$. Keep the pattern: to prove something holds on all of
$[a,b]$, take the supremum of where it holds and show that supremum cannot
be interior. It reappears in the proof of the mean value theorem and in
Chapter 5's appendix.}
\provenance{ECNU \cjk{总练习题}{review} 7 \#2 (bk.\ p.\,160).}

\item Prove that a continuous real function on $[a,b]$ is uniformly
continuous, using the finite covering theorem.\kn{Rudin's 4.19, but assigned
here --- and the point of assigning it here is that Chapter 2's machinery is
already enough. Each $x$ gives $\delta_x$ for $\varepsilon/2$; the balls of
radius $\delta_x/2$ cover $[a,b]$; take a finite subcover and let $\delta$ be
the smallest radius. The halving is what makes two nearby points land in a
common ball. Compare S4.19 in the Chapter 4 appendix, which names the
$\delta$ this construction produces.}
\provenance{ECNU \cjk{习题}{Ex} 7.1 \#11 (bk.\ p.\,155). Companion \#10:
prove the root-existence theorem by the same method.}

\item Let $(x_n)$ be monotone. Prove that if it has a cluster point, that
cluster point is unique and equals $\sup_n x_n$ (for increasing
$(x_n)$).\kn{Monotonicity collapses the whole cluster set to a point --- the
sequence version of the fact that a monotone function has one-sided limits
everywhere. Contrast with $\sin(n\pi/4)$ at S2.2, which has five: order is
exactly what forbids returning.}
\provenance{ECNU \cjk{习题}{Ex} 7.1 \#7 (bk.\ p.\,155).}

\item Prove that the set of cluster points of the closed interval $[a,b]$ is
$[a,b]$ itself.\kn{The concrete instance of ``$[a,b]$ is perfect,'' which
Rudin asserts in passing at 2.44's discussion but never asks anyone to
verify. Both inclusions need a word: every point of $[a,b]$ is approached
from within, and no point outside is.}
\provenance{ECNU \cjk{习题}{Ex} 7.1 \#6 (bk.\ p.\,155).}

\item Sharpen the nested interval theorem: if $a_n < a_{n+1} < b_{n+1} < b_n$
strictly for all $n$ and $b_n - a_n \to 0$, prove the common point $\xi$
satisfies $a_n < \xi < b_n$ strictly for every $n$.\kn{S2.3 gives only
$a_n \le \xi \le b_n$, and several later arguments need the strict version
--- in particular any argument that must place $\xi$ in the \emph{interior}
of every stage. Cheap to prove, easy to assume without noticing.}
\provenance{ECNU \cjk{习题}{Ex} 7.1 \#3 (bk.\ p.\,155).}

\item Construct a divergent sequence $(a_n)$ with the property that for every
$\varepsilon > 0$ there is an $N$ such that $n > N$ implies
$|a_{n+1} - a_n| < \varepsilon$.\kn{The pseudo-Cauchy counterexample:
consecutive terms getting close does \emph{not} make a sequence Cauchy. The
canonical answer is the partial sums of the harmonic series, where
$|a_{n+1}-a_n| = 1/(n+1) \to 0$ while the sums diverge. Rudin states 3.10
with the full $|a_n - a_m|$ and never asks why the weaker condition will not
do --- but every student proposes it eventually.}
\provenance{USTC Ch.\ 1 suppl.\ \#4 (bk.\ p.\,38).}

% ---------------- Group II ----------------
\item $\star$~\emph{(The completion, built by hand.)} Let $(X,d)$ be a metric
space. Call two Cauchy sequences in $X$ equivalent if
$d(x_n,y_n) \to 0$, write $\overline X$ for the set of equivalence classes,
and define $\overline d\bigl([x],[y]\bigr) = \lim d(x_n,y_n)$. Prove:
(a) the relation is an equivalence relation; (b) the limit defining
$\overline d$ exists and is independent of representatives, and
$\overline d$ is a metric; (c) $\overline X$ is complete; (d) $x \mapsto
[(x,x,x,\dots)]$ embeds $X$ isometrically; (e) the image is dense in
$\overline X$.\kn{The construction of $\R$ from $\Q$, done once in
generality --- Rudin gives it only as Exercise 3.24 and only sketches it.
Part (b) is where the work is: you must prove a limit exists before you can
use it as a definition, and that is what the Cauchy hypothesis is for.
Chapter 1's appendix builds $\R$ by cuts instead; comparing the two is the
best available illustration that ``the'' real numbers are characterised by a
property, not a construction.}
\provenance{Tao II, Ex.\ 1.4.8 (bk.\ p.\,21); Yu Pin \cjk{定理}{Thm} 77
(p.\,123) adds the universal property.}

\item Prove that a metric space is compact if and only if it is complete and
totally bounded (S2.9).\kn{For the hard direction, build nested
subsequences: cover by finitely many balls of radius $1$ and keep one
containing infinitely many terms; refine with radius $1/2$; and so on. The
diagonal subsequence is Cauchy, and completeness converts it into a
convergent one. The diagonal trick is the same one Rudin uses at 7.23 for
Arzel\`a--Ascoli --- worth meeting here first, where it is uncluttered.}
\provenance{Tao II, Ex.\ 1.5.10 (bk.\ p.\,26).}

\item In the space $X$ of absolutely summable sequences with
$d(a,b) = \sum_n |a_n - b_n|$, let $e^{(n)}$ be the sequence with $1$ in
place $n$ and $0$ elsewhere. Prove that $\{e^{(n)} : n \in \N\}$ is closed
and bounded, that $X$ is complete, and that the set is nevertheless not
compact.\kn{Sharper than Rudin's Exercise 2.16, where compactness fails
because $\Q$ is incomplete. Here nothing is incomplete: the failure is
\emph{dimension}. Any two of the $e^{(n)}$ are at distance $2$, so no
subsequence is Cauchy and no finite family of balls of radius $1$ covers the
set. This is the picture drawn at S2.9.}
\provenance{Tao II, Ex.\ 1.5.8 with Ex.\ 1.1.15 (bk.\ pp.\,26, 10).}

\item \emph{(Contraction mapping principle.)} Let $(X,d)$ be complete and
$T : X \to X$ satisfy $d(Tx,Ty) \le \gamma\, d(x,y)$ for some fixed
$\gamma \in (0,1)$. Prove $T$ has exactly one fixed point. Then prove the
refinement: if some iterate $T^r$ is a contraction, $T$ still has exactly
one fixed point --- and check that $x \mapsto e^{-x}$ on $[0,\infty)$ is not
a contraction although $f \circ f$ is.\kn{The first theorem in the book that
genuinely \emph{needs} completeness: the iterates form a Cauchy sequence by
comparison with a geometric series, and only completeness produces the
limit. Rudin proves this at 9.23, five chapters later, for the inverse
function theorem; the one-dimensional case is exercise S4.15 in the Chapter
4 appendix.}
\provenance{Yu Pin \cjk{定理}{Thm} 37 (p.\,77) and the following exercises.}

\item Let $(X,d)$ be a compact metric space and $f : X \to X$. Prove:
(a) if $d(f(x),f(y)) = d(x,y)$ for all $x,y$, then $f$ is a bijection;
(b) if $d(f(x),f(y)) \ge d(x,y)$ for all $x,y$, then $f$ is an isometry;
(c) if $d(f(x),f(y)) \le d(x,y)$ for all $x,y$ and $f$ is onto, then $f$ is
an isometry.\kn{Compactness forcing rigidity. Part (b) is the surprising one
--- a map that never decreases distances on a compact space cannot increase
them either. All three run on the same engine: iterate $f$, extract a
convergent subsequence, and use it to show points that looked separated must
be reachable. Nothing in Rudin's Chapter 2 exercises has this flavour.}
\provenance{Yu Pin, winter problem S (p.\,331).}

\item On a product $Y \times Z$ of metric spaces put
$d\bigl((y_1,z_1),(y_2,z_2)\bigr) =
\sqrt{d_Y(y_1,y_2)^2 + d_Z(z_1,z_2)^2}$. Prove this is a metric, that the
projections are continuous, and that $F : X \to Y \times Z$ is continuous if
and only if both components are. Then prove that the invertible $n \times n$
real matrices form an \emph{open} subset of all matrices and that
$A \mapsto A^{-1}$ is continuous on it.\kn{The universal property of the
product, which is what Rudin's 4.10 (a map into $\R^k$ is continuous iff its
components are) is a special case of. The matrix statements are the payoff:
openness follows because the determinant is continuous and $GL_n$ is the
preimage of $\R \setminus \{0\}$, and continuity of inversion from Cramer's
rule. Chapter 9 will use both without proof.}
\provenance{Yu Pin, homework 4 A8--A12 (p.\,108).}

% ---------------- Group III ----------------
\item Prove that $\R^1$ is uncountable using nested intervals rather than a
diagonal argument: (a) given a closed interval $J$ of positive length and any
$x$, find a closed $I \subset J$ of positive length with $x \notin I$;
(b) given a sequence $(x_n)$, build nested closed intervals $I_n$ with
$x_n \notin I_n$; (c) conclude.\kn{Rudin proves uncountability twice --- by
diagonalisation at 2.14 and, far more powerfully, as a corollary of ``perfect
sets are uncountable'' at 2.43. This is a third proof, and the cheapest: it
uses only S2.3. Note that it is really the same idea as 2.43, with the
nested intervals doing what Rudin's shrinking neighbourhoods do.}
\provenance{Yu Pin, homework 1 B7 (p.\,37).}

\item Prove that any family of pairwise disjoint open intervals in $\R^1$ is
countable.\kn{One line if you see it: each interval contains a rational, and
disjointness makes the assignment injective. The technique --- inject an
unknown family into $\Q$ --- is the standard way to prove countability, and
it is exactly the argument behind Rudin's 4.30 (a monotone function has at
most countably many discontinuities) and behind Exercise 2.29.}
\provenance{Yu Pin, homework 1 B11 (p.\,37).}

\item Prove Cantor's theorem in general: for any set $X$ there is no
surjection $X \to 2^X$. Then prove the Schr\"oder--Bernstein theorem: if
there are injections $A \to B$ and $B \to A$, there is a bijection
$A \to B$.\kn{Cantor's proof is Rudin's 2.14 with the decimals removed:
$A = \{x : x \notin f(x)\}$ is missed by every $f(x)$ --- and note its
kinship with Russell's paradox. Schr\"oder--Bernstein is what makes
cardinality an order relation; the slick proof partitions $A$ using
$D_0 = B \setminus f(A)$ and $D_{n+1} = f(g(D_n))$.}
\provenance{Tao I, Thm 8.3.1 and Ex.\ 8.3.2--8.3.3 (bk.\ pp.\,195, 198).}

\item $\star$~Construct a compact topological space containing a sequence
with no convergent subsequence.\kn{The counterexample proving that S2.10 is
a theorem about \emph{metric} spaces and not about topological ones. Tao's
construction well-orders $\R$, takes the union of all countable initial
segments, adjoins a maximum, and gives the result its order topology; every
point below the maximum has only countably many predecessors, which is what
strands the sequence. Rudin never distinguishes the two compactnesses
because he never leaves metric spaces.}
\provenance{Tao II, Ex.\ 2.5.8 with 2.5.5 (bk.\ pp.\,43, 42).}

% ---------------- Group IV ----------------
\item $\star$~\emph{(Furstenberg.)} On $\Z$ call a set open if it is empty or
a union of arithmetic progressions $U_{a,b} = \{ka+b : k \in \Z\}$ with
$a \ge 1$. Prove: (a) this is a topology; (b) no nonempty finite set is
open, so no set with finite complement is closed; (c) each $U_{a,b}$ is both
open and closed; (d) $\Z \setminus \{-1,1\} = \bigcup_{p \text{ prime}}
U_{p,0}$; (e) conclude that there are infinitely many primes.\kn{The best
advertisement in this chapter for topology as a tool rather than a
vocabulary. If there were finitely many primes, the union in (d) would be a
finite union of closed sets, hence closed, so $\{-1,1\}$ would be open ---
contradicting (b). Every step uses only Rudin's 2.23--2.24, on a space with
no metric in sight.}
\provenance{Yu Pin, homework 5 \cjk{娱乐问题}{recreational problem}
(p.\,130).}

\end{enumerate}
